\section{Fermat fourfold results and the universal transfer problem}
\label{app:fermat-related}

The Fermat-variety results discussed here belong to a separate
manuscript \cite{Bha26fermat}; they do not replace the abelian and K3
constructions of this paper. In particular, Theorem B of that
manuscript is the odd-degree fourfold assertion, whereas the principal
unresolved class here is the generic mixed Mumford class.

Li's preprint \cite[Theorem~1.1]{Li26odd}, deposited on 8 September
2026, already states the rational Hodge conjecture for every Fermat
fourfold
\[
 X_m^4=\{x_0^m+\cdots+x_5^m=0\}\subset\PP^5,
 \qquad m\ge3\ \text{odd}.
\]
It therefore precedes the same theorem-level conclusion in
\cite{Bha26fermat}. Li's later preprint
\cite[Theorem~A and Theorem~7.20]{Li26residual} states an exact
classification of six-entry Hodge characters outside a surface
contraction and one exact $5$-standard multiset. It has six primitive
residual orbits, at degrees $33,39,39,45,105,105$.
This differs from the three terminal orbits of the profile descent in
\cite{Bha26fermat}: a tuple that has no lowering move need not be
residual for those two positive predicates. Appendix C of
\cite{Bha26fermat} gives explicit comparisons. The seed constructions
used in these arguments are credited to Jumagulov
\cite[Theorems~1.2--1.4]{Jum26odd}.
We cite both works of Li as preprints for this comparison; their
unbounded conductor arguments are not hypotheses in our geometric
proofs.

There is a relevant caution for the mixed classes in this paper.
Algebraicity of the Hodge classes of individual factors does not, by a
formal K\"unneth argument, give algebraicity of every Hodge class on
their product. A rational Hodge tensor can involve factors that are
not themselves Hodge classes. For an elliptic curve $E$, put
$H=H^1(E,\QQ)$ and let $D^pH$ be the largest rational sub-Hodge
structure contained in $F^pH_{\CC}$. Complex conjugation shows that
$D^1H=0$: a rational sub-Hodge structure contained in $H^{1,0}$ is
also contained in $H^{0,1}$ and hence vanishes. Nevertheless the
K\"unneth component $\Delta_{1,1}$ of the diagonal is a nonzero
rational class of type $(1,1)$ in $H\otimes H$, so
\[
 D^1(H\otimes H)\ne
 \sum_{a+b=1}D^aH\otimes D^bH=0.
\]
To see nonvanishing without a choice of basis, its action on $H^1(E)$
is the identity, by the diagonal correspondence. This is the
diagonal-class failure recorded in \cite[\S1.1]{Li26odd} and
\cite[Remark~3.21]{Bha26fermat}; it identifies a false tensor identity,
rather than disproving all conclusions of the argument in which that
identity occurs.

An actual algebraic correspondence supplies the required transfer.
For smooth projective $X,Y$, a cycle $Z$ on $X$ and an algebraic
correspondence $\Gamma$ of suitable codimension, compatibility of the
cycle class with intersection and proper push-forward gives
\[
 \Gamma_*\operatorname{cl}(Z)=
 \operatorname{cl}\!\left((\operatorname{pr}_Y)_*
       (\operatorname{pr}_X^*Z\cdot\Gamma)\right).
\]
Indeed, apply pullback compatibility, multiplicativity of the cycle
class, and push-forward compatibility successively. Thus an algebraic
source class and an algebraic correspondence do give an algebraic
image. The missing correspondence cannot be replaced by the
algebraicity of classes on a separate collection of varieties.
The cited Fermat results supply neither the cycle for the generic
mixed Mumford class nor a universal correspondence for \Fiii{}.
Those construction problems remain open here.
